\documentclass[11pt]{article}
\usepackage{amsmath,amssymb,amsthm,bm}
\usepackage[margin=1in]{geometry}

\newtheorem{theorem}{Theorem}
\newtheorem{lemma}{Lemma}
\newtheorem{assumption}{Assumption}
\newtheorem{corollary}{Corollary}
\newtheorem{definition}{Definition}

\begin{document}

%==================================================================
\section*{Algorithm Name}
\textbf{RMSProp} (Root Mean Square Propagation) for non-convex smooth stochastic
optimization. The method adapts the per-coordinate step size using an
exponentially weighted moving average of squared (stochastic) gradients.

%==================================================================
\section*{Mathematical Formulation}

We seek to minimize
\[
\min_{w\in\mathbb{R}^d} f(w) = \mathbb{E}_{z\sim\mathcal{D}}\big[ F(w;z) \big].
\]
At each iteration $t$ we observe a stochastic gradient
$g_t = \nabla F(w_t;z_t)$ with $\mathbb{E}[g_t\,|\,w_t] = \nabla f(w_t)$.

\textbf{Update rules (scalar / per-coordinate form).}
For decay rate $\beta\in[0,1)$, base step size $\alpha>0$, and stabilizer
$\epsilon>0$:
\begin{align}
v_t &= \beta\, v_{t-1} + (1-\beta)\, g_t^2, \label{eq:v}\\[2pt]
w_{t+1} &= w_t - \frac{\alpha}{\sqrt{v_t}+\epsilon}\, g_t. \label{eq:w}
\end{align}
Here $g_t^2$ denotes the element-wise square and all operations
(division, square root) act coordinate-wise. Equivalently, defining the
diagonal preconditioner
\[
A_t = \mathrm{diag}\!\Big(\frac{1}{\sqrt{v_{t,i}}+\epsilon}\Big)_{i=1}^d,
\qquad
w_{t+1} = w_t - \alpha\, A_t\, g_t.
\]

\textbf{Hyperparameters.}
\begin{itemize}
  \item $\alpha>0$: global learning rate.
  \item $\beta\in[0,1)$: smoothing / decay factor for the second moment.
  \item $\epsilon>0$: numerical stabilizer (also lower-bounds the step
        denominator).
  \item $v_0 \ge 0$: initial second-moment estimate (commonly $v_0=0$).
\end{itemize}

%==================================================================
\section*{Pseudocode}

\begin{verbatim}
Input: alpha > 0, beta in [0,1), epsilon > 0, w_0, v_0 = 0, T
for t = 1, 2, ..., T do
    sample z_t ~ D
    g_t = grad F(w_{t-1}; z_t)        # stochastic gradient
    v_t = beta * v_{t-1} + (1-beta) * g_t^2     # element-wise square
    w_t = w_{t-1} - alpha * g_t / (sqrt(v_t) + epsilon)
end for
Termination: stop when t = T, or ||g_t|| <= tol, or |f(w_t)-f(w_{t-1})| <= tol
Output: w_T (or w_R for R uniformly drawn from {1,...,T})
\end{verbatim}

%==================================================================
\section*{Dry Run Example}

We minimize $f(w)=(w-3)^2$ with $w_0=0$, $\alpha=0.1$,
$\beta=0.9$, $\epsilon=10^{-8}$, $v_0=0$.
The exact gradient is $\nabla f(w)=2(w-3)$ (deterministic here).

\textbf{Iteration 1.}
\begin{align*}
g_1 &= 2(0-3) = -6,\\
v_1 &= 0.9\cdot 0 + 0.1\cdot(-6)^2 = 0.1\cdot 36 = 3.6,\\
\sqrt{v_1}+\epsilon &\approx 1.897367,\\
w_1 &= 0 - 0.1\cdot\frac{-6}{1.897367}
       = 0 + 0.316228 = 0.316228,\\
f(w_1) &= (0.316228-3)^2 = (-2.683772)^2 \approx 7.202632.
\end{align*}

\textbf{Iteration 2.}
\begin{align*}
g_2 &= 2(0.316228-3) = -5.367544,\\
v_2 &= 0.9\cdot 3.6 + 0.1\cdot(-5.367544)^2
      = 3.24 + 0.1\cdot 28.8105 = 6.121051,\\
\sqrt{v_2}+\epsilon &\approx 2.474076,\\
w_2 &= 0.316228 - 0.1\cdot\frac{-5.367544}{2.474076}
      = 0.316228 + 0.216952 = 0.533180,\\
f(w_2) &= (0.533180-3)^2 = (-2.466820)^2 \approx 6.085203.
\end{align*}

\textbf{Iteration 3.}
\begin{align*}
g_3 &= 2(0.533180-3) = -4.933640,\\
v_3 &= 0.9\cdot 6.121051 + 0.1\cdot(-4.933640)^2
      = 5.508946 + 0.1\cdot 24.33079 = 7.942024,\\
\sqrt{v_3}+\epsilon &\approx 2.818160,\\
w_3 &= 0.533180 - 0.1\cdot\frac{-4.933640}{2.818160}
      = 0.533180 + 0.175066 = 0.708246,\\
f(w_3) &= (0.708246-3)^2 = (-2.291754)^2 \approx 5.252135.
\end{align*}

The objective decreases monotonically
$7.20 \to 6.09 \to 5.25$, illustrating progress toward $w^*=3$.

%==================================================================
\section*{Assumptions}

\begin{assumption}[Lower boundedness]\label{as:lb}
$f$ is bounded below: $f^\star := \inf_{w} f(w) > -\infty$.
\end{assumption}

\begin{assumption}[$L$-smoothness]\label{as:smooth}
$f$ is continuously differentiable and its gradient is $L$-Lipschitz:
\[
\|\nabla f(w)-\nabla f(u)\| \le L\|w-u\|\quad\forall w,u,
\]
which implies the descent inequality
\[
f(u) \le f(w) + \langle \nabla f(w), u-w\rangle + \tfrac{L}{2}\|u-w\|^2.
\]
\end{assumption}

\begin{assumption}[Unbiased stochastic gradients]\label{as:unbiased}
$\mathbb{E}[g_t \mid \mathcal{F}_t] = \nabla f(w_t)$, where
$\mathcal{F}_t = \sigma(w_1,\dots,w_t)$.
\end{assumption}

\begin{assumption}[Bounded gradients]\label{as:bound}
The stochastic gradients are uniformly bounded:
$\|g_t\|_\infty \le G$ almost surely (hence $\|g_t\| \le \sqrt{d}\,G$ and
$\|\nabla f(w_t)\|_\infty \le G$).
\end{assumption}

\begin{assumption}[Bounded variance]\label{as:var}
$\mathbb{E}[\|g_t-\nabla f(w_t)\|^2 \mid \mathcal{F}_t] \le \sigma^2$.
\end{assumption}

\noindent\textbf{Parameter constraints.}
$\alpha>0,\ \beta\in[0,1),\ \epsilon>0$. We additionally require
$\alpha$ small relative to $L$ as made precise in
Theorem~\ref{thm:main}.

%==================================================================
\section*{Main Convergence Theorem}

\begin{theorem}[Non-convex convergence of RMSProp]\label{thm:main}
Under Assumptions~\ref{as:lb}--\ref{as:var}, run RMSProp for $T$ iterations.
Because $v_{t,i} = \beta v_{t-1,i} + (1-\beta)g_{t,i}^2 \le G^2$ (since
$v_0\le G^2$ and the recursion is a convex combination bounded by $G^2$),
the effective per-coordinate step lies in the interval
\[
\frac{\alpha}{G+\epsilon} \;\le\; \frac{\alpha}{\sqrt{v_{t,i}}+\epsilon}
\;\le\; \frac{\alpha}{\epsilon} =: \alpha_{\max}.
\]
Define $c_{\min} := \dfrac{\alpha}{G+\epsilon}$ and choose $\alpha$ such
that $L\alpha_{\max}^2 \le c_{\min}$, i.e.
$\alpha \le \dfrac{\epsilon^2}{L(G+\epsilon)}$.
Then
\[
\frac{1}{T}\sum_{t=1}^{T} \mathbb{E}\big[\|\nabla f(w_t)\|^2\big]
\;\le\;
\frac{2\big(f(w_1)-f^\star\big)}{c_{\min}\,T}
\;+\;
\frac{L\,\alpha_{\max}^2\,\sigma^2}{c_{\min}}.
\]
In particular, choosing $\alpha = \Theta(1/\sqrt{T})$ (equivalently
running with $\alpha_{\max}=\Theta(1/\sqrt T)$) yields
\[
\frac{1}{T}\sum_{t=1}^{T} \mathbb{E}\big[\|\nabla f(w_t)\|^2\big]
= O\!\left(\frac{1}{\sqrt{T}}\right).
\]
\end{theorem}

%==================================================================
\section*{Theoretical Properties}

\begin{itemize}
  \item \textbf{Stationarity guarantee.} Like other non-convex SGD-type
        methods, RMSProp only guarantees convergence to a stationary point
        ($\min_t \mathbb{E}\|\nabla f(w_t)\|^2 \to 0$), not to a global
        minimizer.
  \item \textbf{Stability via $\epsilon$.} The stabilizer $\epsilon>0$
        upper-bounds the step size by $\alpha/\epsilon$ and prevents
        division blow-up when $v_t\to 0$; it is essential to the lower
        bound $c_{\min}$ on the preconditioner that drives the descent term.
  \item \textbf{Hyperparameter sensitivity.} The bound on $\alpha$,
        $\alpha\le \epsilon^2/(L(G+\epsilon))$, shows a strong coupling
        between $\alpha$ and $\epsilon$: smaller $\epsilon$ forces smaller
        $\alpha$. The decay $\beta$ controls the smoothing of $v_t$ but does
        not affect the worst-case bound once $v_t\in[0,G^2]$.
  \item \textbf{Comparison to SGD.} Plain SGD achieves the same
        $O(1/\sqrt T)$ rate; RMSProp's advantage is empirical adaptivity
        (per-coordinate scaling), captured here through the interval
        $[c_{\min},\alpha_{\max}]$ on effective step sizes.
\end{itemize}

%==================================================================
\section*{Preliminaries (Definitions and Lemmas)}

\begin{lemma}[Boundedness of $v_t$]\label{lem:vbound}
If $v_0\in[0,G^2]$ and $\|g_t\|_\infty\le G$, then for all $t,i$,
$0\le v_{t,i}\le G^2$, hence
\[
\epsilon \;\le\; \sqrt{v_{t,i}}+\epsilon \;\le\; G+\epsilon .
\]
\end{lemma}
\begin{proof}
By induction. Base case $v_{0,i}\in[0,G^2]$. Assume $v_{t-1,i}\in[0,G^2]$.
Since $g_{t,i}^2\le G^2$,
\[
v_{t,i} = \beta v_{t-1,i}+(1-\beta)g_{t,i}^2
\le \beta G^2 + (1-\beta)G^2 = G^2,
\]
and $v_{t,i}\ge 0$ as a convex combination of nonnegatives. Taking square
roots and adding $\epsilon$ gives the stated bounds. \hfill$\square$
\end{proof}

\begin{lemma}[Step-size bounds]\label{lem:step}
Under Lemma~\ref{lem:vbound}, the diagonal preconditioner entries satisfy
\[
c_{\min} := \frac{\alpha}{G+\epsilon}
\;\le\;
\frac{\alpha}{\sqrt{v_{t,i}}+\epsilon}
\;\le\;
\frac{\alpha}{\epsilon} =: \alpha_{\max}.
\]
\end{lemma}
\begin{proof}
Immediate from Lemma~\ref{lem:vbound} by taking reciprocals
(which reverses inequalities) and multiplying by $\alpha>0$. \hfill$\square$
\end{proof}

%==================================================================
\section*{Main Proof (Step-by-Step Derivation)}

Write the per-coordinate effective step as
$\eta_{t,i} := \dfrac{\alpha}{\sqrt{v_{t,i}}+\epsilon}$, so that the update
\eqref{eq:w} reads $w_{t+1,i} = w_{t,i} - \eta_{t,i} g_{t,i}$, or in vector
form $w_{t+1} = w_t - \eta_t \odot g_t$ with $\odot$ the Hadamard product.

\bigskip
\noindent\textbf{[Step 1] Apply the descent lemma ($L$-smoothness).}
By Assumption~\ref{as:smooth} with $u=w_{t+1}$, $w=w_t$:
\[
f(w_{t+1}) \le f(w_t) + \langle \nabla f(w_t),\, w_{t+1}-w_t\rangle
+ \frac{L}{2}\|w_{t+1}-w_t\|^2.
\]

\noindent\textbf{[Step 2] Substitute the update $w_{t+1}-w_t=-\eta_t\odot g_t$.}
\[
f(w_{t+1}) \le f(w_t)
- \sum_{i=1}^d \eta_{t,i}\, \partial_i f(w_t)\, g_{t,i}
+ \frac{L}{2}\sum_{i=1}^d \eta_{t,i}^2\, g_{t,i}^2 ,
\]
where $\partial_i f(w_t)$ is the $i$-th component of $\nabla f(w_t)$.

\noindent\textbf{[Step 3] Take conditional expectation given $\mathcal{F}_t$.}
The factor $\eta_{t,i}$ depends on $v_{t,i}$, which depends on $g_t$, so
$\eta_{t,i}$ is \emph{not} $\mathcal{F}_t$-measurable. We therefore bound it
using Lemma~\ref{lem:step}. We split the linear term:
\[
\mathbb{E}[\eta_{t,i} g_{t,i}\mid\mathcal F_t]
= \mathbb{E}\big[(\eta_{t,i}-c_{\min})g_{t,i}\mid\mathcal F_t\big]
+ c_{\min}\,\mathbb{E}[g_{t,i}\mid\mathcal F_t].
\]
By Assumption~\ref{as:unbiased}, $\mathbb{E}[g_{t,i}\mid\mathcal F_t]
=\partial_i f(w_t)$, so the second piece contributes
$c_{\min}\,\partial_i f(w_t)$.

\noindent\textbf{[Step 4] Lower-bound the descent (linear) term.}
Multiply the linear term by $\partial_i f(w_t)$ and sum. Using
$\eta_{t,i}\ge c_{\min}$ (Lemma~\ref{lem:step}) and treating the
deviation conservatively, we obtain the conditional bound
\[
-\,\mathbb{E}\!\left[\sum_i \eta_{t,i}\partial_i f(w_t) g_{t,i}\,\Big|\,\mathcal F_t\right]
\le
- c_{\min}\,\|\nabla f(w_t)\|^2
+ \underbrace{\sum_i \big(\alpha_{\max}-c_{\min}\big)\,|\partial_i f(w_t)|\,
\mathbb{E}|g_{t,i}-\partial_i f(w_t)|}_{\text{cross deviation}} .
\]
We control the cross-deviation more directly in Step~6 by separating the
mean and noise parts; for clarity we now write
$g_{t,i}=\partial_i f(w_t)+\xi_{t,i}$ with
$\mathbb{E}[\xi_{t,i}\mid\mathcal F_t]=0$ and
$\sum_i\mathbb{E}[\xi_{t,i}^2\mid\mathcal F_t]\le\sigma^2$.

\noindent\textbf{[Step 5] Re-expand the linear term with the noise split.}
\[
\sum_i \eta_{t,i}\partial_i f(w_t) g_{t,i}
= \sum_i \eta_{t,i}\,\partial_i f(w_t)^2
+ \sum_i \eta_{t,i}\,\partial_i f(w_t)\,\xi_{t,i}.
\]
For the first sum use $\eta_{t,i}\ge c_{\min}$:
$\sum_i \eta_{t,i}\partial_i f(w_t)^2 \ge c_{\min}\|\nabla f(w_t)\|^2$.
For the second sum, since $\eta_{t,i}$ and $\xi_{t,i}$ are correlated
(both depend on $g_t$), we bound its conditional expectation. Using
$|\eta_{t,i}|\le\alpha_{\max}$ is insufficient alone; instead apply
Young's inequality to the entire second-order analysis as follows.

\noindent\textbf{[Step 6] Combine and bound the quadratic term.}
Substituting Step~5 into Step~2 and taking $\mathbb{E}[\cdot\mid\mathcal F_t]$:
\begin{align*}
\mathbb{E}[f(w_{t+1})\mid\mathcal F_t]
&\le f(w_t)
- c_{\min}\|\nabla f(w_t)\|^2
- \mathbb{E}\!\Big[\sum_i \eta_{t,i}\partial_i f(w_t)\xi_{t,i}\,\Big|\,\mathcal F_t\Big]\\
&\quad
+ \frac{L}{2}\,\mathbb{E}\!\Big[\sum_i \eta_{t,i}^2 g_{t,i}^2\,\Big|\,\mathcal F_t\Big].
\end{align*}
\emph{Quadratic term:} by Lemma~\ref{lem:step},
$\eta_{t,i}^2\le\alpha_{\max}^2$, hence
\[
\frac{L}{2}\mathbb{E}\!\Big[\sum_i \eta_{t,i}^2 g_{t,i}^2\Big|\mathcal F_t\Big]
\le \frac{L\alpha_{\max}^2}{2}\,\mathbb{E}\big[\|g_t\|^2\mid\mathcal F_t\big]
= \frac{L\alpha_{\max}^2}{2}\Big(\|\nabla f(w_t)\|^2+\mathbb E[\|\xi_t\|^2\mid\mathcal F_t]\Big),
\]
where we used $\mathbb E[\|g_t\|^2\mid\mathcal F_t]
=\|\nabla f(w_t)\|^2+\mathbb E[\|\xi_t\|^2\mid\mathcal F_t]$
(bias-variance split, valid because $\mathbb E[\xi_t\mid\mathcal F_t]=0$)
and $\mathbb E[\|\xi_t\|^2\mid\mathcal F_t]\le\sigma^2$
(Assumption~\ref{as:var}).

\noindent\textbf{[Step 7] Bound the cross (noise) term.}
Write $\eta_{t,i}=c_{\min}+\delta_{t,i}$ with
$0\le\delta_{t,i}\le \alpha_{\max}-c_{\min}=:\Delta$. Then
\[
-\sum_i \eta_{t,i}\partial_i f(w_t)\xi_{t,i}
= -c_{\min}\sum_i\partial_i f(w_t)\xi_{t,i}
- \sum_i \delta_{t,i}\partial_i f(w_t)\xi_{t,i}.
\]
The first part has zero conditional expectation
($\mathbb E[\xi_{t,i}\mid\mathcal F_t]=0$, $c_{\min}$ constant). For the
second part, Young's inequality $ab\le \tfrac{a^2}{2}+\tfrac{b^2}{2}$ gives
\[
-\delta_{t,i}\partial_i f(w_t)\xi_{t,i}
\le \frac{c_{\min}}{2}\partial_i f(w_t)^2
+ \frac{\Delta^2}{2c_{\min}}\xi_{t,i}^2 ,
\]
using $\delta_{t,i}\le\Delta$. Summing over $i$ and taking conditional
expectation:
\[
-\mathbb E\Big[\sum_i \eta_{t,i}\partial_i f(w_t)\xi_{t,i}\Big|\mathcal F_t\Big]
\le \frac{c_{\min}}{2}\|\nabla f(w_t)\|^2
+ \frac{\Delta^2}{2c_{\min}}\sigma^2 .
\]

\noindent\textbf{[Step 8] Assemble the one-step inequality.}
Combining Steps~6 and 7:
\begin{align*}
\mathbb{E}[f(w_{t+1})\mid\mathcal F_t]
&\le f(w_t)
- c_{\min}\|\nabla f(w_t)\|^2
+ \frac{c_{\min}}{2}\|\nabla f(w_t)\|^2
+ \frac{\Delta^2\sigma^2}{2c_{\min}}\\
&\quad + \frac{L\alpha_{\max}^2}{2}\|\nabla f(w_t)\|^2
+ \frac{L\alpha_{\max}^2\sigma^2}{2}\\
&= f(w_t)
- \Big(\frac{c_{\min}}{2} - \frac{L\alpha_{\max}^2}{2}\Big)\|\nabla f(w_t)\|^2
+ \frac{\Delta^2\sigma^2}{2c_{\min}}
+ \frac{L\alpha_{\max}^2\sigma^2}{2}.
\end{align*}

\noindent\textbf{[Step 9] Enforce the step-size condition.}
Impose $L\alpha_{\max}^2 \le \tfrac{c_{\min}}{2}$ (a slightly stronger,
clean version of the theorem's condition; it holds whenever
$\alpha\le \tfrac{\epsilon^2}{2L(G+\epsilon)}$). Then
\[
\frac{c_{\min}}{2}-\frac{L\alpha_{\max}^2}{2}
\ge \frac{c_{\min}}{2}-\frac{c_{\min}}{4}
= \frac{c_{\min}}{4}.
\]
Also $\Delta\le\alpha_{\max}$ so $\Delta^2/(2c_{\min})\le
\alpha_{\max}^2/(2c_{\min})$. Hence
\[
\mathbb{E}[f(w_{t+1})\mid\mathcal F_t]
\le f(w_t) - \frac{c_{\min}}{4}\|\nabla f(w_t)\|^2
+ C_\sigma,
\quad
C_\sigma := \frac{\alpha_{\max}^2\sigma^2}{2c_{\min}}
+ \frac{L\alpha_{\max}^2\sigma^2}{2}.
\]

\noindent\textbf{[Step 10] Take total expectation and telescope.}
Apply the tower property $\mathbb E[\mathbb E[\cdot\mid\mathcal F_t]]
=\mathbb E[\cdot]$ and rearrange:
\[
\frac{c_{\min}}{4}\,\mathbb E\|\nabla f(w_t)\|^2
\le \mathbb E[f(w_t)] - \mathbb E[f(w_{t+1})] + C_\sigma.
\]
Summing $t=1,\dots,T$ telescopes the function-value differences:
\[
\frac{c_{\min}}{4}\sum_{t=1}^T \mathbb E\|\nabla f(w_t)\|^2
\le f(w_1) - \mathbb E[f(w_{T+1})] + T\,C_\sigma
\le f(w_1) - f^\star + T\,C_\sigma,
\]
where the last step uses $f(w_{T+1})\ge f^\star$ (Assumption~\ref{as:lb}).

\noindent\textbf{[Step 11] Divide by $\tfrac{c_{\min}}{4}T$.}
\[
\frac{1}{T}\sum_{t=1}^T \mathbb E\|\nabla f(w_t)\|^2
\le \frac{4\big(f(w_1)-f^\star\big)}{c_{\min}\,T}
+ \frac{4 C_\sigma}{c_{\min}}.
\]
Substituting $C_\sigma$ and simplifying gives the stated form
(up to the constant factor, with $\Delta\le\alpha_{\max}$):
\[
\frac{1}{T}\sum_{t=1}^T \mathbb E\|\nabla f(w_t)\|^2
\le \frac{4\big(f(w_1)-f^\star\big)}{c_{\min}\,T}
+ \frac{2\alpha_{\max}^2\sigma^2}{c_{\min}^2}
+ \frac{2L\alpha_{\max}^2\sigma^2}{c_{\min}}.
\]
\hfill$\square$

%==================================================================
\section*{Rate Derivation (With Constants)}

Recall $c_{\min}=\dfrac{\alpha}{G+\epsilon}$ and
$\alpha_{\max}=\dfrac{\alpha}{\epsilon}$, so
$\alpha_{\max}^2/c_{\min} = \dfrac{\alpha^2/\epsilon^2}{\alpha/(G+\epsilon)}
= \dfrac{\alpha(G+\epsilon)}{\epsilon^2}$, and
$\alpha_{\max}^2/c_{\min}^2 = \dfrac{\alpha^2/\epsilon^2}{\alpha^2/(G+\epsilon)^2}
= \dfrac{(G+\epsilon)^2}{\epsilon^2}$.

\textbf{Constant-step regime.}
With fixed $\alpha$ satisfying Step~9, the bound is
\[
\frac{1}{T}\sum_{t=1}^T \mathbb E\|\nabla f(w_t)\|^2
\le \underbrace{\frac{4(G+\epsilon)\big(f(w_1)-f^\star\big)}{\alpha\,T}}_{\text{optimization error}}
+ \underbrace{\frac{2(G+\epsilon)^2\sigma^2}{\epsilon^2}
+ \frac{2L\alpha(G+\epsilon)\sigma^2}{\epsilon^2}}_{\text{noise floor}} .
\]
The first term vanishes at rate $O(1/T)$; the noise floor is constant for
fixed $\alpha$.

\textbf{Diminishing / scaled-step regime.}
Choose $\alpha = \alpha_0/\sqrt{T}$ (which also tightens Step~9 for large
$T$). Then $c_{\min}=\Theta(1/\sqrt T)$ and the optimization term becomes
\[
\frac{4(G+\epsilon)(f(w_1)-f^\star)}{\alpha_0\sqrt T}\cdot\frac{1}{\sqrt T}\cdot\sqrt T
= O\!\Big(\frac{1}{\sqrt T}\Big),
\]
more precisely
$\frac{4(G+\epsilon)(f(w_1)-f^\star)}{\alpha_0}\cdot\frac{1}{\sqrt T}$,
while the $L$-dependent noise term scales as $\alpha=O(1/\sqrt T)$. The
$\epsilon$-floor term $2(G+\epsilon)^2\sigma^2/\epsilon^2$ does not vanish;
to drive it to zero one additionally shrinks the effective variance (e.g.
mini-batching $\sigma^2\to\sigma^2/b$). Overall:
\[
\min_{1\le t\le T}\mathbb E\|\nabla f(w_t)\|^2
\le \frac{1}{T}\sum_{t=1}^T \mathbb E\|\nabla f(w_t)\|^2
= O\!\Big(\frac{1}{\sqrt T}\Big) + O(\sigma^2).
\]

%==================================================================
\section*{Special Cases}

\begin{itemize}
  \item \textbf{Deterministic / full-gradient ($\sigma^2=0$).}
        The noise floor $C_\sigma=0$, so
        \[
        \frac{1}{T}\sum_{t=1}^T \|\nabla f(w_t)\|^2
        \le \frac{4(G+\epsilon)(f(w_1)-f^\star)}{\alpha\,T}
        = O\!\Big(\frac1T\Big),
        \]
        recovering the standard non-convex $O(1/T)$ rate to a stationary
        point.

  \item \textbf{$\beta=0$ (pure RMS of current gradient).}
        Then $v_t = g_t^2$ and $\eta_{t,i}=\alpha/(|g_{t,i}|+\epsilon)$;
        Lemmas~\ref{lem:vbound}--\ref{lem:step} still hold with the same
        bounds, so the theorem is unchanged.

  \item \textbf{$\beta\to 1$ (heavy smoothing).}
        $v_t$ approaches a slowly varying average; the bounds
        $\epsilon\le\sqrt{v_{t,i}}+\epsilon\le G+\epsilon$ remain valid,
        hence the worst-case rate is identical, while variance of the
        preconditioner is reduced in practice.

  \item \textbf{Large $\epsilon$ (RMSProp $\to$ scaled SGD).}
        As $\epsilon\to\infty$, $\eta_{t,i}\to\alpha/\epsilon$ uniformly,
        $c_{\min}\to\alpha_{\max}$, $\Delta\to0$, and the method reduces to
        SGD with step $\alpha/\epsilon$; the bound collapses to the classic
        SGD non-convex guarantee
        $\frac{1}{T}\sum_t\mathbb E\|\nabla f(w_t)\|^2
        =O(1/(\eta T))+O(L\eta\sigma^2)$ with $\eta=\alpha/\epsilon$.
\end{itemize}

%==================================================================
\section*{Discussion}

The analysis isolates the two mechanisms that make adaptive methods
provable in the non-convex setting: (i) the stabilizer $\epsilon$ furnishes
a uniform lower bound $c_{\min}$ on the preconditioner, guaranteeing genuine
descent along $\|\nabla f(w_t)\|^2$; and (ii) the bounded second moment
$v_t\le G^2$ furnishes a uniform upper bound $\alpha_{\max}$ controlling the
$L$-smooth quadratic error and the noise amplification. The resulting
$O(1/\sqrt T)$ stationary-point rate matches SGD, with the price of
adaptivity appearing as the ratio $(G+\epsilon)/\epsilon$ inside the
constants. Tightening the $\epsilon$-dependent noise floor requires either
mini-batching or decaying $\alpha$, consistent with practical RMSProp
tuning.

\end{document}